% Part: lambda-calculus
% Chapter: introduction
% Section: minimization

\documentclass[../../../include/open-logic-section]{subfiles}

\begin{document}

\olfileid{lam}{int}{min}
\olsection{\usetoken{S}{lambda definable} تابعې د کمينه‌موندنې تر عمل لاندې تړلې دي}

\begin{lem}
فرض کړئ $f(x,y)$ !!{lambda definable} ده۔ $g$ داسې تعريف کړئ:
\[
g(x) \simeq \umin{y}{f(x,y)}.
\]
نو $g$ هم !!{lambda definable} ده۔
\end{lem}

\begin{proof}
مفکوره په لنډ ډول داسې ده۔ د $x$ په لرلو سره به د ثابت ټکي لامبډا
ترم $Y$ په وسيله داسې تابع $h_x(n)$ تعريف کړو چې له~$n$ څخه د~$y$
لټون پيلوي؛ بيا $g(x)$ يوازې $h_x(0)$ دے۔ تابع~$h_x$ د ثابت ټکي
د يوې معادلې د حل په توګه ښودلے شو:
\[
h_x(n) \simeq
\begin{cases}
n & \text{که $f(x,n) = 0$ وي} \\
h_x(n+1) & \text{کنه.}
\end{cases}
\]

تفصيل ئې دا دے۔ څرنګه چې $f$ لامبډا-تعريفېدونکې ده، يو ترم
$F$ ئې !!{lambda defined} کوي۔
په ياد ولرئ چې يو لامبډا ترم~$D$ هم لرو، داسې چې
$D(M, N, \bar{0}) \red M$ او $D(M, N, \bar{1})
\red N$۔ د اوس دپاره $x$ ثابت ونيسئ۔ د $h_x$ د !!{lambda define}
کولو دپاره داسې ترم~$H$ غواړو چې پر~$x$ تړلے وي او دا معادله پوره کړي:
\[
H(\num{n}) \equiv D(\num{n}, H(S(\num{n})), F(x, \num{n})).
\]
دا کار د ثابت ټکي ترم~$Y$ په وسيله کولے شو۔ لومړے $U$ داسې ترم
ونيسئ:
\[
\lambd[h][\lambd[z][D(z,(h(Sz)),F(x,z))]],
\]
او بيا $H$ ترم~$YU$ ونيسئ۔ پام وکړئ چې په~$H$ کښې يوازنے آزاد
متغير~$x$ دے۔ راځئ وښيو چې $H$ پاسنۍ معادله پوره کوي۔

د $Y$ له تعريفه لرو
\[
H = YU \equiv U(YU) = U(H).
\]
په ځانګړې توګه، د هر طبيعي عدد~$n$ دپاره لرو
\begin{align*}
H(\num{n}) & \equiv U(H, \num{n}) \\
& \red D(\num{n}, H(S(\num{n})), F(x, \num{n})),
\end{align*}
لکه څنګه چې غوښتل شوي وو۔ پام وکړئ چې که په وروستۍ کرښه کښې د~$x$
په ځاے عددنښه $\num{m}$ کښېږدئ، نو که $F(\num{m}, \num{n})$
$\num{0}$ ته تقليل شي، دا عبارت $\num{n}$ ته تقليلېږي؛ او که
$F(\num{m}, \num{n})$ بلې هرې عددنښې ته تقليل شي، نو عبارت
$H(S(\num{n}))$ ته تقليلېږي۔

په تعريف شوي حالت کښې $G$ ترم $\lambd[x][H(\num 0)]$ ونيسئ۔ د هر~$m$ دپاره که $g(m)$ تعريف شوې وي، له صفره تر شاهد پورې ټولې ازموينې تعريف شوې دي۔ نو $G(\num m)$، $\overline {g(m)}$ ته تقليلېږي۔ په دې حالت کښې $G$،~$g$ !!{lambda define}s، يعنې غوښتل شوې عددنښه ورکوي۔
د ناتعريفه ننوت لپاره د دې ټاکونکي مستقيمه بڼه په عمومي توګه بس نۀ ده۔ يو بې‌عادي‌بڼې ترم کولے شي په تطبيق کښې هغه برخه وغورځوي چې عادي بڼه نۀ لري۔ د عمومي لمې لپاره د \olref[cmp]{thm:lambda-computable} له مخې د ازموينې جزوي محاسبه واخلو۔ پر صفر يې پيلوو، او تر هغه پر بل عدد ازموينه کوو چې اوسنی ارزښت تعريف شوے او غيرصفر وي۔ د لومړي صفر پر موندلو هماغه شاخص ورکوو۔ که کوم اړين ارزښت ناتعريف وي، يا تل غيرصفر ارزښتونه راشي، لټون نۀ درېږي۔ دا د جزوي کمينه‌موندنې تعريف دے۔ د \olref[lrp]{thm:computable-lambda} له مخې د دغې جزوي محاسبې لپاره يو لامبډا ترم شته چې پر ناتعريفه ننوت عادي بڼه نۀ لري۔
\end{proof}

\end{document}
